\section{Statistical certificate and experimental protocol}
\label{sec:protocol}

The theorem gives a probability bound; an experiment needs a decision rule. The rule below separates three ingredients that are often conflated: sampling uncertainty, measured bias in each branch, and uncertainty in whether the laboratory really implements the theoretical null.

\subsection{Fixed-sample certificate}

Let $(K_{\mathrm A},N_{\mathrm A})$ and $(K_{\mathrm R},N_{\mathrm R})$ be successes and fixed trials in the two branches. Define
\[
\widehat S=\lambda\frac{K_{\mathrm A}}{N_{\mathrm A}}
+(1-\lambda)\frac{K_{\mathrm R}}{N_{\mathrm R}}.
\]
For independent Bernoulli trials, the one-sided weighted Hoeffding radius is
\begin{equation}
r_\alpha=
\sqrt{\frac12\left(
\frac{\lambda^2}{N_{\mathrm A}}+
\frac{(1-\lambda)^2}{N_{\mathrm R}}
\right)\log\frac1\alpha}.
\label{eq:hoeffding}
\end{equation}
Before observing outcomes, let $\delta_{\mathrm A}$ and $\delta_{\mathrm R}$ upper-bound possible upward bias in the two estimators, and let $\delta_{\mathrm{null}}$ enlarge the theoretical null for source, isolation, and interface uncertainty. We certify incompatibility with the declared classical-memory null only if
\begin{equation}
\widehat S-\lambda\delta_{\mathrm A}
-(1-\lambda)\delta_{\mathrm R}-r_\alpha
>\Cstream+\delta_{\mathrm{null}}.
\label{eq:certificate}
\end{equation}
Under independent fixed-sample trials, the false-positive probability is at most $\alpha$. Choosing $n$ or $\lambda$ after seeing data, optional stopping, or selecting the best of several lengths requires a different analysis.

For an alternative score $S_{\mathrm{alt}}$, define
\[
g=S_{\mathrm{alt}}-\lambda\delta_{\mathrm A}
-(1-\lambda)\delta_{\mathrm R}
-\Cstream-\delta_{\mathrm{null}}.
\]
If $g\leq0$, no amount of data certifies the claim under those allowances. If $g>0$, fixed sample sizes satisfying $r_\alpha+r_\beta<g$ guarantee power at least $1-\beta$ by a second Hoeffding bound. The implementation allocates trials asymptotically as $N_{\mathrm A}:N_{\mathrm R}=\lambda:(1-\lambda)$.

\subsection{Commitment and branch timing}

A valid trial follows the left-to-right order of \cref{fig:protocol}. The verifier first prepares and independently characterises the EPR sources. It releases $A_i$ only after $B_{i-1}$ has entered trusted sequestration and, after receiving $B_n$, requires delivery of the terminal memory port $M$. Only then is the hidden AUDIT/RETURN choice generated. The preregistered branch decoder and binary acceptance measurement are applied, and every loss or exclusion is recorded under a predeclared rule. The branch may be randomised per trial, but the final branch counts and stopping rule must be fixed or analysed with a sequentially valid method.

\subsection{Mandatory shortcut controls}

These controls are part of the scientific claim, not optional diagnostics. Deliberately dephasing the persistent carrier should remove the coherent advantage, whereas making a strong classical record should improve AUDIT while lowering global RETURN. Removing sequestration should recover the collective curve in \eqref{eq:collective}; allowing the device to revisit an earlier $B_i$ should likewise expose why the causal boundary matters. Decoder controls replace the coherent inverse by identity or by a mismatched inverse. Independent source and measurement calibration must bound false EPR acceptance and parity bias, and a blind-choice control must test whether branch information leaks before commitment. Finally, a side-channel audit must bound coherent transfer through buses, pulse tails, clocks, resonators, unused levels, and shared ancillas.

The null does not forbid inaccessible erasure environments. Consequently, visible reset of $M$ is a reproducibility and wiring check, not a certificate that every physical copy of a classical transcript has vanished.

\subsection{Reproducible analysis interface}

The public command-line tool evaluates the exact bound, plans trial counts, and applies the fixed-sample rule through three strict-JSON commands:
\begin{small}
\begin{verbatim}
carmenq bound --steps 8
carmenq plan --steps 8 --forecast-model
carmenq analyse --steps 8 \
  --audit-successes 9700 --audit-trials 10000 \
  --return-successes 9500 --return-trials 10000
\end{verbatim}
\end{small}
Systematic allowances, $\alpha$, $\beta$, and null slack are command-line inputs and appear in the output. An example preregistration records the terminal port, source calibration, decoder, exclusions, controls, and stopping rule.
